\documentclass[11pt]{article}
\usepackage[a4paper,top=0.86in,bottom=0.92in,left=0.82in,right=0.82in,headheight=22pt]{geometry}
\usepackage[T1]{fontenc}
\usepackage[utf8]{inputenc}
\usepackage{lmodern}
\usepackage{microtype}
\usepackage{setspace}
\usepackage{parskip}
\usepackage{xcolor}
\usepackage{colortbl}
\usepackage{graphicx}
\usepackage{amsmath,amssymb}
\usepackage{fancyvrb}
\usepackage{booktabs}
\usepackage{longtable}
\usepackage{array}
\usepackage{calc}
\usepackage{tabularx}
\usepackage{makecell}
\usepackage{multirow}
\usepackage{adjustbox}
\usepackage{caption}
\usepackage{float}
\usepackage{enumitem}
\usepackage{titlesec}
\usepackage{fancyhdr}
\usepackage{lastpage}
\usepackage{hyperref}
\usepackage{xurl}
\usepackage{etoolbox}
\usepackage{pdflscape}
\usepackage{ragged2e}
\usepackage{seqsplit}

\definecolor{reportBlack}{HTML}{111111}
\definecolor{reportGray}{HTML}{666666}
\definecolor{ruleGray}{HTML}{BDBDBD}
\definecolor{tableHead}{HTML}{F2F4F7}
\definecolor{codeBg}{HTML}{F7F7F7}
\definecolor{linkBlue}{HTML}{1F4E79}

\hypersetup{colorlinks=true, linkcolor=linkBlue, urlcolor=linkBlue, citecolor=linkBlue}
\setcounter{secnumdepth}{-\maxdimen}
\urlstyle{same}
\setlength{\emergencystretch}{4em}
\tolerance=2000
\hbadness=2000
\sloppy
\hyphenation{Con-fi-gu-ra-tion Im-ple-men-ta-tion Eval-u-a-tion Re-pro-duc-i-bil-i-ty}
\setstretch{1.05}
\setlist[itemize]{leftmargin=1.2em,itemsep=2pt,topsep=2pt}
\setlist[enumerate]{leftmargin=1.4em,itemsep=2pt,topsep=2pt}
\captionsetup{labelfont=bf,font=small}

\pagestyle{fancy}
\fancyhf{}
\fancyhead[L]{\small\color{reportGray}CSI500 Stock Selection --- Project Report}
\fancyhead[R]{\small\color{reportGray}ML Competition \textperiodcentered{} Spring 2026}
\fancyfoot[C]{\small\color{reportGray}\thepage}
\renewcommand{\headrulewidth}{0.4pt}
\renewcommand{\headrule}{\hbox to\headwidth{\color{ruleGray}\leaders\hrule height \headrulewidth\hfill}}
\renewcommand{\footrulewidth}{0.4pt}
\renewcommand{\footrule}{\hbox to\headwidth{\color{ruleGray}\leaders\hrule height \footrulewidth\hfill}}

\titleformat{\section}{\Large\bfseries\color{reportBlack}}{\thesection}{0.6em}{}
\titleformat{\subsection}{\large\bfseries\color{reportBlack}}{\thesubsection}{0.6em}{}
\titleformat{\subsubsection}{\normalsize\bfseries\color{reportBlack}}{\thesubsubsection}{0.6em}{}
\titlespacing*{\section}{0pt}{1.6em}{0.6em}
\titlespacing*{\subsection}{0pt}{1.1em}{0.4em}
\titlespacing*{\subsubsection}{0pt}{0.8em}{0.25em}

\renewcommand{\arraystretch}{1.18}
\setlength{\extrarowheight}{2pt}
\setlength{\tabcolsep}{4.5pt}
\setlength{\arrayrulewidth}{0.35pt}
\arrayrulecolor{ruleGray}
\AtBeginEnvironment{longtable}{%
  \footnotesize
  \setlength{\tabcolsep}{2.4pt}%
  \renewcommand{\arraystretch}{1.30}%
}
\AfterEndEnvironment{longtable}{%
  \normalsize
  \setlength{\tabcolsep}{4.5pt}%
  \renewcommand{\arraystretch}{1.18}%
}
\newcolumntype{Y}{>{\RaggedRight\arraybackslash}X}
\newcolumntype{Z}{>{\Centering\arraybackslash}X}
\newcolumntype{R}{>{\Centering\arraybackslash}X}
\newcommand{\code}[1]{{\ttfamily\footnotesize\seqsplit{#1}}}
\providecommand{\tightlist}{\setlength{\itemsep}{0pt}\setlength{\parskip}{0pt}}

\usepackage{color}
\usepackage{fancyvrb}
\newcommand{\VerbBar}{|}
\newcommand{\VERB}{\Verb[commandchars=\\\{\}]}
\DefineVerbatimEnvironment{Highlighting}{Verbatim}{commandchars=\\\{\}}
% Add ',fontsize=\small' for more characters per line
\newenvironment{Shaded}{}{}
\newcommand{\AlertTok}[1]{\textcolor[rgb]{1.00,0.00,0.00}{\textbf{#1}}}
\newcommand{\AnnotationTok}[1]{\textcolor[rgb]{0.38,0.63,0.69}{\textbf{\textit{#1}}}}
\newcommand{\AttributeTok}[1]{\textcolor[rgb]{0.49,0.56,0.16}{#1}}
\newcommand{\BaseNTok}[1]{\textcolor[rgb]{0.25,0.63,0.44}{#1}}
\newcommand{\BuiltInTok}[1]{\textcolor[rgb]{0.00,0.50,0.00}{#1}}
\newcommand{\CharTok}[1]{\textcolor[rgb]{0.25,0.44,0.63}{#1}}
\newcommand{\CommentTok}[1]{\textcolor[rgb]{0.38,0.63,0.69}{\textit{#1}}}
\newcommand{\CommentVarTok}[1]{\textcolor[rgb]{0.38,0.63,0.69}{\textbf{\textit{#1}}}}
\newcommand{\ConstantTok}[1]{\textcolor[rgb]{0.53,0.00,0.00}{#1}}
\newcommand{\ControlFlowTok}[1]{\textcolor[rgb]{0.00,0.44,0.13}{\textbf{#1}}}
\newcommand{\DataTypeTok}[1]{\textcolor[rgb]{0.56,0.13,0.00}{#1}}
\newcommand{\DecValTok}[1]{\textcolor[rgb]{0.25,0.63,0.44}{#1}}
\newcommand{\DocumentationTok}[1]{\textcolor[rgb]{0.73,0.13,0.13}{\textit{#1}}}
\newcommand{\ErrorTok}[1]{\textcolor[rgb]{1.00,0.00,0.00}{\textbf{#1}}}
\newcommand{\ExtensionTok}[1]{#1}
\newcommand{\FloatTok}[1]{\textcolor[rgb]{0.25,0.63,0.44}{#1}}
\newcommand{\FunctionTok}[1]{\textcolor[rgb]{0.02,0.16,0.49}{#1}}
\newcommand{\ImportTok}[1]{\textcolor[rgb]{0.00,0.50,0.00}{\textbf{#1}}}
\newcommand{\InformationTok}[1]{\textcolor[rgb]{0.38,0.63,0.69}{\textbf{\textit{#1}}}}
\newcommand{\KeywordTok}[1]{\textcolor[rgb]{0.00,0.44,0.13}{\textbf{#1}}}
\newcommand{\NormalTok}[1]{#1}
\newcommand{\OperatorTok}[1]{\textcolor[rgb]{0.40,0.40,0.40}{#1}}
\newcommand{\OtherTok}[1]{\textcolor[rgb]{0.00,0.44,0.13}{#1}}
\newcommand{\PreprocessorTok}[1]{\textcolor[rgb]{0.74,0.48,0.00}{#1}}
\newcommand{\RegionMarkerTok}[1]{#1}
\newcommand{\SpecialCharTok}[1]{\textcolor[rgb]{0.25,0.44,0.63}{#1}}
\newcommand{\SpecialStringTok}[1]{\textcolor[rgb]{0.73,0.40,0.53}{#1}}
\newcommand{\StringTok}[1]{\textcolor[rgb]{0.25,0.44,0.63}{#1}}
\newcommand{\VariableTok}[1]{\textcolor[rgb]{0.10,0.09,0.49}{#1}}
\newcommand{\VerbatimStringTok}[1]{\textcolor[rgb]{0.25,0.44,0.63}{#1}}
\newcommand{\WarningTok}[1]{\textcolor[rgb]{0.38,0.63,0.69}{\textbf{\textit{#1}}}}

\newcommand{\studentName}{Yueyao Chen, yc7009}
\newcommand{\reportTitle}{CSI500 Stock Selection}
\newcommand{\reportSubtitle}{Final Report}

\begin{document}

\begin{titlepage}
\thispagestyle{fancy}
\centering
\vspace*{\fill}
{\Huge\bfseries \reportTitle\par}
\vspace{0.35cm}
{\LARGE\color{reportGray}\reportSubtitle\par}
\vspace{0.9cm}
{\Large \studentName\par}
\vspace*{\fill}
\end{titlepage}

\tableofcontents
\newpage

\textbf{Project:} CSI500 Stock Selection ML Competition, Spring 2026\\
\textbf{Final Stage2 portfolio:} \code{stage2\_final\_portfolio.csv}\\
\textbf{Final as-of date:} 2026-05-08\\
\textbf{Final production entry point:}
\code{stage2\_baseline\_guard\_ensemble.py}\\
\textbf{Local data coverage:} 2023-01-03 to 2026-05-08

\section{Executive Summary}\label{executive-summary}

The task is to submit a long-only CSI500 portfolio for a future holding
window. The submission file contains only two columns,
\code{stock\_code} and \code{weight}; the scoring script then
computes the realized portfolio return over the evaluation window and
subtracts the CSI500 benchmark return. Therefore, the modeling problem
has two separate layers:

\begin{enumerate}
\def\labelenumi{\arabic{enumi}.}
\tightlist
\item
  Predict a cross-sectional ordering or confidence score for stocks
  using only information available at the submission date.
\item
  Convert those scores into a valid portfolio with at least 30
  positive-weight stocks, non-negative weights, weights summing to one,
  and no single stock above the 10\% cap.
\end{enumerate}

The final Stage2 route is a \textbf{Regime-gated ensemble}. It contains
several candidate experts, including an official-style XGBoost fallback,
a LightGBM/XGBoost tree consensus, a weekly alpha overlay, a full-week
cycle tree, and meta/hybrid fallback routes. A market-state gate chooses
the route using only as-of observable index and breadth features. For
the final as-of date, 2026-05-08, the gate selected the defensive
official-style XGBoost branch because the market looked like a broad
overheated rebound.

On the final strict self-test of 12 complete Monday-Friday five-day
windows, the final adaptive route achieved:

\begin{longtable}[]{@{}|
  >{\RaggedRight\arraybackslash}m{(\linewidth - 12\tabcolsep - 7\arrayrulewidth) * \real{0.1304}}|
  >{\Centering\arraybackslash}m{(\linewidth - 12\tabcolsep - 7\arrayrulewidth) * \real{0.1739}}|
  >{\Centering\arraybackslash}m{(\linewidth - 12\tabcolsep - 7\arrayrulewidth) * \real{0.1739}}|
  >{\Centering\arraybackslash}m{(\linewidth - 12\tabcolsep - 7\arrayrulewidth) * \real{0.1739}}|
  >{\Centering\arraybackslash}m{(\linewidth - 12\tabcolsep - 7\arrayrulewidth) * \real{0.1739}}|
  >{\Centering\arraybackslash}m{(\linewidth - 12\tabcolsep - 7\arrayrulewidth) * \real{0.1739}}|@{}}
\hline\rowcolor{tableHead}
\begin{minipage}[c]{\linewidth}\centering\bfseries
Model
\end{minipage} & \begin{minipage}[c]{\linewidth}\centering\bfseries
Mean excess
\end{minipage} & \begin{minipage}[c]{\linewidth}\centering\bfseries
Median excess
\end{minipage} & \begin{minipage}[c]{\linewidth}\centering\bfseries
Min excess
\end{minipage} & \begin{minipage}[c]{\linewidth}\centering\bfseries
Max excess
\end{minipage} & \begin{minipage}[c]{\linewidth}\centering\bfseries
Negative windows
\end{minipage} \\
\hline
\endhead
\hline
\endlastfoot
Final regime-gated ensemble & +4.001\% & +3.001\% & +0.907\% & +12.308\%
& 0 / 12 \\
\hline
Original XGBoost baseline & +0.600\% & +0.264\% & -1.221\% & +2.916\% &
5 / 12 \\
\end{longtable}

The final route improved both the average excess return and the downside
floor relative to the original baseline. The final report intentionally
presents failed attempts as well, because many high-looking experiments
were later rejected after stricter window definition and leakage checks.

The logic of the final method is:

\begin{longtable}[]{@{}|
  >{\RaggedRight\arraybackslash}m{(\linewidth - 6\tabcolsep - 4\arrayrulewidth) * \real{0.3333}}|
  >{\RaggedRight\arraybackslash}m{(\linewidth - 6\tabcolsep - 4\arrayrulewidth) * \real{0.3333}}|
  >{\RaggedRight\arraybackslash}m{(\linewidth - 6\tabcolsep - 4\arrayrulewidth) * \real{0.3333}}|@{}}
\hline\rowcolor{tableHead}
\begin{minipage}[c]{\linewidth}\centering\bfseries
Step
\end{minipage} & \begin{minipage}[c]{\linewidth}\centering\bfseries
Decision
\end{minipage} & \begin{minipage}[c]{\linewidth}\centering\bfseries
Why it comes next
\end{minipage} \\
\hline
\endhead
\hline
\endlastfoot
Scoring analysis & Model outputs weights, while
\code{score\_submission.py} computes realized excess return &
Separates prediction from evaluation \\
\hline
Factor design & Build lagged price, liquidity, risk, market-relative,
and weekly-cycle signals & Defines what the model is allowed to know at
as-of \\
\hline
Model search & Test XGBoost, LightGBM, CatBoost, LSTM, direct
rank/weight learning, and hybrids & Identifies which learners are stable
versus regime-sensitive \\
\hline
Validation correction & Switch from generic five-trading-day windows to
complete Monday-Friday self-test windows & Aligns the model-selection
test with the Stage2 objective \\
\hline
Final route & Use an as-of regime gate to choose among experts & Avoids
forcing one model into every market state \\
\hline
Portfolio shape & Choose top-k and weight concentration conditional on
the selected route & Converts scores into a compliant, risk-controlled
submission \\
\end{longtable}

\section{1. Factors}\label{factors}

\subsection{1.1 Prediction Target And Scoring
Definition}\label{prediction-target-and-scoring-definition}

The official scoring script applies portfolio weights at the open of the
evaluation start date and holds through the close of the evaluation end
date. For stock (i), the realized return is:

\begin{Shaded}
\begin{Highlighting}[]
\NormalTok{r\_i = close\_i(end) / close\_i(day\_before\_start) {-} 1}
\end{Highlighting}
\end{Shaded}

The portfolio return is the weighted sum of realized stock returns:

\begin{Shaded}
\begin{Highlighting}[]
\NormalTok{portfolio\_return = sum\_i weight\_i * r\_i}
\end{Highlighting}
\end{Shaded}

The official competition metric is excess return:

\begin{Shaded}
\begin{Highlighting}[]
\NormalTok{excess\_return = portfolio\_return {-} CSI500\_benchmark\_return}
\end{Highlighting}
\end{Shaded}

This means the model does not need to output excess return directly. It
outputs stock weights. The excess return is computed only after future
data become known in \code{score\_submission.py}.

During development, we tested three target designs:

\begin{longtable}[]{@{}|
  >{\RaggedRight\arraybackslash}m{(\linewidth - 8\tabcolsep - 5\arrayrulewidth) * \real{0.2500}}|
  >{\RaggedRight\arraybackslash}m{(\linewidth - 8\tabcolsep - 5\arrayrulewidth) * \real{0.2500}}|
  >{\RaggedRight\arraybackslash}m{(\linewidth - 8\tabcolsep - 5\arrayrulewidth) * \real{0.2500}}|
  >{\RaggedRight\arraybackslash}m{(\linewidth - 8\tabcolsep - 5\arrayrulewidth) * \real{0.2500}}|@{}}
\hline\rowcolor{tableHead}
\begin{minipage}[c]{\linewidth}\centering\bfseries
Target design
\end{minipage} & \begin{minipage}[c]{\linewidth}\centering\bfseries
Implementation idea
\end{minipage} & \begin{minipage}[c]{\linewidth}\centering\bfseries
Result
\end{minipage} & \begin{minipage}[c]{\linewidth}\centering\bfseries
Final decision
\end{minipage} \\
\hline
\endhead
\hline
\endlastfoot
Raw future return & Predict \code{target\_5d} and rank stocks by
predicted return & Stable, close to official baseline, useful fallback &
Retained \\
\hline
Future excess return & Predict
\code{target\_excess\_5d\ =\ stock\ future\ return\ -\ index\ future\ return}
& Useful for weekly cycle experts & Retained in specialized route \\
\hline
Direct future rank/weight & Train toward future rank or synthetic
optimal weight & Unstable; some portfolios became too equal-like or too
noisy & Rejected as primary route \\
\end{longtable}

The final route separates \textbf{Prediction} from \textbf{Portfolio
construction}. This was more robust than directly learning weights,
because the ``true weight'' is not an observed label in the data; it is
an artificial transformation of future returns and can amplify noise.

\subsection{1.2 Official Data Used}\label{official-data-used}

The final production route uses only the official local data cache.
External data were explored, cleaned, and tested, but not promoted into
the final route because they did not improve both mean excess and
downside stability.

\begin{longtable}[]{@{}|>{\RaggedRight\arraybackslash}m{0.30\linewidth}|>{\RaggedRight\arraybackslash}m{0.42\linewidth}|>{\Centering\arraybackslash}m{0.18\linewidth}|@{}}
\hline\rowcolor{tableHead}
\multicolumn{1}{|c|}{\textbf{Data item}} & \multicolumn{1}{c|}{\textbf{File}} & \multicolumn{1}{c|}{\textbf{Value}} \\
\hline
\endhead
\hline
\endlastfoot
Stock daily panel & \code{data/prices.parquet} & 397,757 rows \\
\hline
Unique stocks & \code{data/prices.parquet} & 499 \\
\hline
Trading dates & \code{data/prices.parquet} & 807 \\
\hline
Date range & \code{data/prices.parquet} & 2023-01-03 to 2026-05-08 \\
\hline
Index daily panel & \code{data/index.parquet} & 807 rows \\
\hline
Constituent list & \code{data/constituents.csv} & 499 rows \\
\end{longtable}

The core stock data contain daily OHLCV-style fields such as open, high,
low, close, volume, amount, and turnover. The index data provide CSI500
close and return features used for benchmark-relative signals.

\subsection{1.3 Core Feature Set}\label{core-feature-set}

The production baseline path uses a compact set of technical and
cross-sectional features defined in \code{features.py} as
\code{CORE\_FEATURE\_COLUMNS}.

\begin{longtable}[]{@{}|
  >{\RaggedRight\arraybackslash}m{(\linewidth - 6\tabcolsep - 4\arrayrulewidth) * \real{0.3333}}|
  >{\RaggedRight\arraybackslash}m{(\linewidth - 6\tabcolsep - 4\arrayrulewidth) * \real{0.3333}}|
  >{\RaggedRight\arraybackslash}m{(\linewidth - 6\tabcolsep - 4\arrayrulewidth) * \real{0.3333}}|@{}}
\hline\rowcolor{tableHead}
\begin{minipage}[c]{\linewidth}\centering\bfseries
Feature group
\end{minipage} & \begin{minipage}[c]{\linewidth}\centering\bfseries
Features
\end{minipage} & \begin{minipage}[c]{\linewidth}\centering\bfseries
Rationale
\end{minipage} \\
\hline
\endhead
\hline
\endlastfoot
Short and medium returns & \code{ret\_1d}, \code{ret\_5d},
\code{ret\_10d}, \code{ret\_20d}, \code{ret\_60d} & Capture
momentum, reversal, and trend persistence \\
\hline
Volatility and risk & \code{vol\_20d} & Avoid stocks whose recent risk
dominates signal \\
\hline
Liquidity and activity & \code{volume\_z\_20d},
\code{turnover\_ma\_20d} & Identify abnormal trading activity and
tradeability \\
\hline
Trend location & \code{close\_over\_ma20}, \code{close\_over\_ma60},
\code{rsi\_14} & Measure overextension versus moving averages \\
\hline
Cross-sectional ranks & \code{ret\_5d\_rank}, \code{ret\_20d\_rank},
\code{vol\_20d\_rank} & Normalize signals across a trading day \\
\end{longtable}

The compact core was retained because blindly adding features repeatedly
improved one or two windows while hurting the out-of-sample floor. Tree
models can use both raw and ranked features; mechanically dropping
highly correlated raw/rank pairs reduced performance in several route
checks.

\subsection{1.4 Expanded Feature Families
Tested}\label{expanded-feature-families-tested}

Several richer feature groups were implemented and used by specialized
experts or rejected after ablation.

\begin{longtable}[]{@{}|
  >{\RaggedRight\arraybackslash}m{(\linewidth - 8\tabcolsep - 5\arrayrulewidth) * \real{0.2500}}|
  >{\RaggedRight\arraybackslash}m{(\linewidth - 8\tabcolsep - 5\arrayrulewidth) * \real{0.2500}}|
  >{\RaggedRight\arraybackslash}m{(\linewidth - 8\tabcolsep - 5\arrayrulewidth) * \real{0.2500}}|
  >{\RaggedRight\arraybackslash}m{(\linewidth - 8\tabcolsep - 5\arrayrulewidth) * \real{0.2500}}|@{}}
\hline\rowcolor{tableHead}
\begin{minipage}[c]{\linewidth}\centering\bfseries
Feature family
\end{minipage} & \begin{minipage}[c]{\linewidth}\centering\bfseries
Examples
\end{minipage} & \begin{minipage}[c]{\linewidth}\centering\bfseries
Why tested
\end{minipage} & \begin{minipage}[c]{\linewidth}\centering\bfseries
Final use
\end{minipage} \\
\hline
\endhead
\hline
\endlastfoot
Momentum and reversal & \code{ret\_3d}, \code{mom\_accel\_5\_20},
\code{ret\_60d\_rank} & Balance short-term bursts with medium trend &
Used in tree/weekly experts \\
\hline
Risk and drawdown & \code{vol\_ratio\_5\_20},
\code{downside\_vol\_20d}, \code{drawdown\_20d} & Penalize unstable
names and recent crash risk & Used in defensive routes \\
\hline
Market-relative features & \code{excess\_ret\_5d},
\code{excess\_ret\_20d}, \code{beta\_60d},
\code{residual\_ret\_20d} & Separate stock alpha from market beta &
Used in weekly and quality routes \\
\hline
Flow and price-volume features & \code{amount\_z\_20d},
\code{obv\_20d}, \code{price\_volume\_corr\_20d} & Approximate
money-flow confirmation & Used in weekly alpha overlay \\
\hline
Intraday and gap features & \code{gap\_mean\_20d},
\code{intraday\_mean\_20d}, \code{overnight\_vol\_20d} & Detect
reopening risk and persistent demand & Used in weekly alpha
experiments \\
\hline
Calendar cycle features & Monday/Friday, month start/end, post-gap,
full-workweek flags & Align Stage2 with five-day workweek behavior &
Used in weekly cycle route \\
\hline
External/open data & Valuation, regime, fund-flow fields & Potentially
add non-price information & Rejected from final \\
\hline
Graph/peer features & Correlation-neighbor return features & Model group
influence among stocks & Rejected from final \\
\end{longtable}

\subsection{1.5 Weekly Cycle Features}\label{weekly-cycle-features}

After we realized that the final Stage2 target was better represented by
complete five-day workweeks, we added known-in-advance calendar cycle
features. These do not use future prices. They are functions of the
as-of date and the next five business days:

\begin{longtable}[]{@{}|
  >{\RaggedRight\arraybackslash}m{(\linewidth - 4\tabcolsep - 3\arrayrulewidth) * \real{0.5000}}|
  >{\RaggedRight\arraybackslash}m{(\linewidth - 4\tabcolsep - 3\arrayrulewidth) * \real{0.5000}}|@{}}
\hline\rowcolor{tableHead}
\begin{minipage}[c]{\linewidth}\centering\bfseries
Weekly feature
\end{minipage} & \begin{minipage}[c]{\linewidth}\centering\bfseries
Interpretation
\end{minipage} \\
\hline
\endhead
\hline
\endlastfoot
\code{eval\_starts\_monday} & Whether the holding window begins on
Monday \\
\hline
\code{eval\_ends\_friday} & Whether the holding window ends on
Friday \\
\hline
\code{eval\_is\_full\_workweek} & Whether the evaluation is a complete
Monday-Friday week \\
\hline
\code{weekend\_gap\_days} & Length of the non-trading gap before the
holding window \\
\hline
\code{eval\_month\_start} and \code{eval\_month\_end} & Month-start
or month-end flow regime \\
\hline
\code{weekly\_risk\_appetite} & Interaction of Monday risk appetite
with recent returns and OBV \\
\hline
\code{weekly\_friday\_derisk} & Friday de-risking tilt using trend
quality and volatility \\
\hline
\code{month\_start\_flow} & Month-start liquidity and flow
confirmation \\
\hline
\code{month\_end\_defensive} & Month-end defensive tilt \\
\hline
\code{post\_gap\_reopen\_flow} & Reopening behavior after weekend or
holiday gaps \\
\hline
\code{weekly\_carry\_quality} & Full-week carry quality using OBV,
trend, close location, and volatility \\
\end{longtable}

These features were used only in the weekly expert route. They were not
forced into the compact baseline fallback, because the fallback
performed best as a simpler and more regularized model.

\subsection{1.6 Feature Normalization And Correlation
Handling}\label{feature-normalization-and-correlation-handling}

The final feature-processing approach was deliberately moderate:

\begin{longtable}[]{@{}|
  >{\RaggedRight\arraybackslash}m{(\linewidth - 6\tabcolsep - 4\arrayrulewidth) * \real{0.3333}}|
  >{\RaggedRight\arraybackslash}m{(\linewidth - 6\tabcolsep - 4\arrayrulewidth) * \real{0.3333}}|
  >{\RaggedRight\arraybackslash}m{(\linewidth - 6\tabcolsep - 4\arrayrulewidth) * \real{0.3333}}|@{}}
\hline\rowcolor{tableHead}
\begin{minipage}[c]{\linewidth}\centering\bfseries
Processing step
\end{minipage} & \begin{minipage}[c]{\linewidth}\centering\bfseries
Implementation
\end{minipage} & \begin{minipage}[c]{\linewidth}\centering\bfseries
Reason
\end{minipage} \\
\hline
\endhead
\hline
\endlastfoot
Rolling features & Computed only from current and historical
observations per stock & Avoid future leakage \\
\hline
Cross-sectional rank features & Per-date rank percentile for selected
signals & Reduce scale drift across market regimes \\
\hline
Robust normalization in weekly route & Rank-normalize stock-level
features within each date & Stabilize tree learners \\
\hline
Correlation filtering in weekly route & Spearman correlation filter,
threshold 0.90 & Remove redundant route-specific features \\
\hline
Target clipping in weekly route & Clip training target at 1\% and 99\%
quantiles & Reduce outlier dominance \\
\hline
Time decay & Half-life 180 days in weekly route & Emphasize recent
regimes without discarding history \\
\end{longtable}

We did not perform aggressive de-correlation globally. In this project,
high feature correlation did not always mean a feature was useless for
tree splits. The safer compromise was to use compact features for the
baseline route and correlation filtering only in richer specialized
routes.

\section{2. Models}\label{models}

\subsection{2.1 Model Development
Timeline}\label{model-development-timeline}

The final route came from a long sequence of model and feature
experiments. The important part is that the project did not move
linearly; some promising paths were later rejected after better
validation.

\begin{longtable}[]{@{}|
  >{\RaggedRight\arraybackslash}m{(\linewidth - 8\tabcolsep - 5\arrayrulewidth) * \real{0.2500}}|
  >{\RaggedRight\arraybackslash}m{(\linewidth - 8\tabcolsep - 5\arrayrulewidth) * \real{0.2500}}|
  >{\RaggedRight\arraybackslash}m{(\linewidth - 8\tabcolsep - 5\arrayrulewidth) * \real{0.2500}}|
  >{\RaggedRight\arraybackslash}m{(\linewidth - 8\tabcolsep - 5\arrayrulewidth) * \real{0.2500}}|@{}}
\hline\rowcolor{tableHead}
\begin{minipage}[c]{\linewidth}\centering\bfseries
Phase
\end{minipage} & \begin{minipage}[c]{\linewidth}\centering\bfseries
Experiment
\end{minipage} & \begin{minipage}[c]{\linewidth}\centering\bfseries
Observation
\end{minipage} & \begin{minipage}[c]{\linewidth}\centering\bfseries
Decision
\end{minipage} \\
\hline
\endhead
\hline
\endlastfoot
Baseline understanding & Read \code{score\_submission.py} and official
XGBoost baseline & Model predicts scores/weights; score script computes
excess return & Keep prediction and scoring separate \\
\hline
Official XGBoost & \code{baseline\_xgboost.py}, raw 5-day target,
top-50 rank weights & Simple and stable, but weak mean excess & Keep as
reference and fallback \\
\hline
Direct rank/weight learning & Target future rank and synthetic target
weight & Weight labels were unstable and sometimes produced near-equal
weights & Reject as main route \\
\hline
Feature pruning & Test core OHLCV, ranks, volatility, liquidity, trend,
RSI & Compact feature set had better floor than broad stuffing & Keep
compact core \\
\hline
LSTM sequence model & Learn short-term sequences and confidence & Strong
on several early rolling windows & Not stable across horizon changes \\
\hline
Layered/gated LSTM & LSTM plus LightGBM confirmation and confidence
weights & Useful for Stage1 three-day objective & Not final for
Stage2 \\
\hline
LightGBM/XGBoost consensus & Blend two tree portfolio rankings &
Improved some regimes & Retained as expert ingredient \\
\hline
CatBoost robust challenger & Quantile/MAE-style robust model & Not
competitive enough standalone & Keep as rejected comparison \\
\hline
Regularized consensus & Stricter tree ensemble and risk overlays &
Helped in narrow regimes & Retained as fallback route \\
\hline
Open-data enhancement & Strict lagging/cleaning of valuation, regime,
fund flow & Noisy, did not improve mean and min together & Reject from
final \\
\hline
Full-week correction & Change validation to complete Monday-Friday
windows & Invalidated old optimistic results & Reset Stage2 standard \\
\hline
Weekly alpha route & Flow, OBV, low-vol confidence, calendar cycle &
Useful in selected regimes & Retained as expert \\
\hline
Final regime gate & Route selection by as-of market state & Best strict
12-window mean and floor & Final route \\
\end{longtable}

\subsection{2.2 Official-Style XGBoost
Baseline}\label{official-style-xgboost-baseline}

The baseline model predicts raw five-day future stock return, then ranks
the cross-section.

\begin{longtable}[]{@{}|>{\RaggedRight\arraybackslash}m{0.58\linewidth}|>{\Centering\arraybackslash}m{0.32\linewidth}|@{}}
\hline\rowcolor{tableHead}
\multicolumn{1}{|c|}{\textbf{Setting}} & \multicolumn{1}{c|}{\textbf{Value}} \\
\hline
\endhead
\hline
\endlastfoot
Target & \code{target\_5d} \\
\hline
estimators & 400 \\
\hline
Max depth & 5 \\
\hline
Learning rate & 0.05 \\
\hline
subsample & 0.8 \\
\hline
Colsample by tree & 0.8 \\
\hline
Min child weight & 10 \\
\hline
L2 regularization & 1.0 \\
\hline
Tree method & \code{hist} \\
\hline
Validation days & 10 trading days \\
\hline
Embargo & 5 trading days \\
\hline
Early stopping & 30 rounds \\
\hline
Default portfolio size & 50 names \\
\end{longtable}

The split for each as-of date is time ordered:

\begin{Shaded}
\begin{Highlighting}[]
\NormalTok{available data up to as\_of}
\NormalTok{train target cutoff = as\_of {-} 5 trading days}

\NormalTok{[ historical train rows ] [ 5{-}day embargo ] [ 10{-}day validation rows ] [ as\_of prediction ]}
\end{Highlighting}
\end{Shaded}

The embargo is important because the target for date (t) uses prices
through (t+5). Without the embargo, training labels could overlap with
validation prices.

\subsection{2.3 LightGBM / XGBoost Tree
Consensus}\label{lightgbm-xgboost-tree-consensus}

The tree consensus route was inspired by the idea that different tree
learners can agree on robust names even when raw predicted scores
differ. The route trains a tuned XGBoost and a LightGBM-style portfolio
helper, then combines their selected stocks and ranks.

\begin{longtable}[]{@{}|
  >{\RaggedRight\arraybackslash}m{(\linewidth - 4\tabcolsep - 3\arrayrulewidth) * \real{0.5000}}|
  >{\RaggedRight\arraybackslash}m{(\linewidth - 4\tabcolsep - 3\arrayrulewidth) * \real{0.5000}}|@{}}
\hline\rowcolor{tableHead}
\begin{minipage}[c]{\linewidth}\centering\bfseries
Component
\end{minipage} & \begin{minipage}[c]{\linewidth}\centering\bfseries
Role
\end{minipage} \\
\hline
\endhead
\hline
\endlastfoot
Tuned XGBoost & Strong tabular ranking model \\
\hline
LightGBM & Fast gradient boosting challenger \\
\hline
Consensus rank & Prefer names selected by both models \\
\hline
Rank-weight layer & Avoid pathological raw-score scaling \\
\hline
Optional factor reweight & Apply drawdown/vol/quality tilts only when
regime gate allows \\
\end{longtable}

The most useful finding was that model averaging alone was not
sufficient. The consensus needed a portfolio layer that preserved
conviction while controlling concentration.

\subsection{2.4 Weekly Cycle Tree
Expert}\label{weekly-cycle-tree-expert}

The weekly cycle route trains on all available five-trading-day targets
but adds features describing the next five-business-day calendar shape.
It is not trained only on full weeks, which helps avoid overfitting to
too few examples.

\begin{longtable}[]{@{}|
  >{\RaggedRight\arraybackslash}m{(\linewidth - 8\tabcolsep - 5\arrayrulewidth) * \real{0.2000}}|
  >{\Centering\arraybackslash}m{(\linewidth - 8\tabcolsep - 5\arrayrulewidth) * \real{0.2667}}|
  >{\Centering\arraybackslash}m{(\linewidth - 8\tabcolsep - 5\arrayrulewidth) * \real{0.2667}}|
  >{\Centering\arraybackslash}m{(\linewidth - 8\tabcolsep - 5\arrayrulewidth) * \real{0.2667}}|@{}}
\hline\rowcolor{tableHead}
\begin{minipage}[c]{\linewidth}\centering\bfseries
Setting
\end{minipage} & \begin{minipage}[c]{\linewidth}\centering\bfseries
LightGBM
\end{minipage} & \begin{minipage}[c]{\linewidth}\centering\bfseries
XGBoost
\end{minipage} & \begin{minipage}[c]{\linewidth}\centering\bfseries
CatBoost challenger
\end{minipage} \\
\hline
\endhead
\hline
\endlastfoot
Robust loss & Huber & Pseudo-Huber & Quantile \\
\hline
Estimators/iterations & 650 & 520 & 520 \\
\hline
Learning rate & 0.022 & 0.022 & 0.025 \\
\hline
Depth/leaves & 47 leaves & Depth 3 & Depth 5 \\
\hline
Min child / leaf & 70 & 45 & L2 leaf reg 10 \\
\hline
subsample & 0.82 & 0.82 & Bagging temp 0.35 \\
\hline
colsample & 0.70 & 0.70 & Ordered boosting style \\
\hline
regularization & Alpha 0.35, lambda 5.0 & Alpha 0.20, lambda 6.0 &
Random strength 1.0 \\
\end{longtable}

The final weekly cycle configuration used the LightGBM/XGBoost pair, not
CatBoost, because the robust two-tree pair gave better stability in the
final full-week checks.

\subsection{2.5 LSTM And Confidence-Weight
Experiments}\label{lstm-and-confidence-weight-experiments}

The LSTM route was developed because stock data are sequential. The main
idea was to learn a short-term sequence representation, then combine it
with a confidence score:

\begin{Shaded}
\begin{Highlighting}[]
\NormalTok{final confidence = model agreement + rank strength {-} volatility/drawdown penalty}
\end{Highlighting}
\end{Shaded}

For Stage1, this was useful. On the three-day validation around the
Stage1 horizon, the best LSTM/LightGBM confidence configuration reached
about +3.3\% mean excess on three recent windows:

\begin{longtable}[]{@{}|
  >{\RaggedRight\arraybackslash}m{(\linewidth - 16\tabcolsep - 9\arrayrulewidth) * \real{0.0968}}|
  >{\Centering\arraybackslash}m{(\linewidth - 16\tabcolsep - 9\arrayrulewidth) * \real{0.1290}}|
  >{\Centering\arraybackslash}m{(\linewidth - 16\tabcolsep - 9\arrayrulewidth) * \real{0.1290}}|
  >{\Centering\arraybackslash}m{(\linewidth - 16\tabcolsep - 9\arrayrulewidth) * \real{0.1290}}|
  >{\Centering\arraybackslash}m{(\linewidth - 16\tabcolsep - 9\arrayrulewidth) * \real{0.1290}}|
  >{\Centering\arraybackslash}m{(\linewidth - 16\tabcolsep - 9\arrayrulewidth) * \real{0.1290}}|
  >{\Centering\arraybackslash}m{(\linewidth - 16\tabcolsep - 9\arrayrulewidth) * \real{0.1290}}|
  >{\Centering\arraybackslash}m{(\linewidth - 16\tabcolsep - 9\arrayrulewidth) * \real{0.1290}}|@{}}
\hline\rowcolor{tableHead}
\begin{minipage}[c]{\linewidth}\centering\bfseries
Config.
\end{minipage} & \begin{minipage}[c]{\linewidth}\centering\bfseries
Alpha\_lstm
\end{minipage} & \begin{minipage}[c]{\linewidth}\centering\bfseries
Risk penalty
\end{minipage} & \begin{minipage}[c]{\linewidth}\centering\bfseries
Top\_k
\end{minipage} & \begin{minipage}[c]{\linewidth}\centering\bfseries
Rank power
\end{minipage} & \begin{minipage}[c]{\linewidth}\centering\bfseries
Confidence power
\end{minipage} & \begin{minipage}[c]{\linewidth}\centering\bfseries
Mean excess
\end{minipage} & \begin{minipage}[c]{\linewidth}\centering\bfseries
Min excess
\end{minipage} \\
\hline
\endhead
\hline
\endlastfoot
Best recent Stage1 LSTM confidence & 0.95 & 0.30 & 40 & 1.35 & 3.0 &
+3.302\% & +3.246\% \\
\end{longtable}

However, this did not transfer reliably to Stage2. Once the target
changed to strict five-day workweeks, LSTM routes became more
regime-sensitive and could produce negative windows. The final Stage2
route therefore keeps LSTM in the historical fallback chain but does not
use it as the final primary expert.

\subsection{2.6 Why Direct Weight Learning Was
Rejected}\label{why-direct-weight-learning-was-rejected}

Directly learning future rank and target weight looked attractive
because the submission itself is a weight file. In practice, it caused
two problems:

\begin{longtable}[]{@{}|
  >{\RaggedRight\arraybackslash}m{(\linewidth - 4\tabcolsep - 3\arrayrulewidth) * \real{0.5000}}|
  >{\RaggedRight\arraybackslash}m{(\linewidth - 4\tabcolsep - 3\arrayrulewidth) * \real{0.5000}}|@{}}
\hline\rowcolor{tableHead}
\begin{minipage}[c]{\linewidth}\centering\bfseries
Issue
\end{minipage} & \begin{minipage}[c]{\linewidth}\centering\bfseries
Why it happened
\end{minipage} \\
\hline
\endhead
\hline
\endlastfoot
Synthetic labels are noisy & There is no true observed ``best weight'';
any label is constructed from future returns \\
\hline
Weight scale instability & Small return differences can create very
different target weights \\
\hline
Equal-looking portfolios & Regularization and noisy labels sometimes
collapsed weights toward similar values \\
\hline
Poor transfer across horizons & 3-Day, generic 5-day, and full-week
5-day objectives favored different shapes \\
\end{longtable}

The final design uses score/rank learning plus a separate portfolio
allocator. This made the model easier to validate and easier to
constrain under competition rules.

\subsection{2.7 Final Regime-Gated
Ensemble}\label{final-regime-gated-ensemble}

The final production route is implemented in
\code{stage2\_baseline\_guard\_ensemble.py}. It is a
mixture-of-experts selector, not a simple average of all model outputs.

\begin{longtable}[]{@{}|
  >{\RaggedRight\arraybackslash}m{(\linewidth - 6\tabcolsep - 4\arrayrulewidth) * \real{0.3333}}|
  >{\RaggedRight\arraybackslash}m{(\linewidth - 6\tabcolsep - 4\arrayrulewidth) * \real{0.3333}}|
  >{\RaggedRight\arraybackslash}m{(\linewidth - 6\tabcolsep - 4\arrayrulewidth) * \real{0.3333}}|@{}}
\hline\rowcolor{tableHead}
\begin{minipage}[c]{\linewidth}\centering\bfseries
Expert
\end{minipage} & \begin{minipage}[c]{\linewidth}\centering\bfseries
Implementation
\end{minipage} & \begin{minipage}[c]{\linewidth}\centering\bfseries
Role
\end{minipage} \\
\hline
\endhead
\hline
\endlastfoot
Official-style XGBoost fallback & \code{baseline\_xgboost.py} and
fallback function inside final route & Stable defensive branch \\
\hline
LightGBM/XGBoost tree consensus & \code{stage2\_tree\_consensus.py}
plus \code{lightgbm\_portfolio.py} and
\code{tuned\_xgboost\_portfolio.py} & Tabular consensus branch \\
\hline
Weekly alpha overlay & \code{stage2\_weekly\_alpha\_overlay.py} &
Flow/volatility/OBV style confidence branch \\
\hline
Weekly consensus & \code{stage2\_weekly\_consensus\_ensemble.py} &
Combines weekly alpha and weekly cycle ideas \\
\hline
Weekly cycle tree & \code{stage2\_weekly\_cycle\_tree.py} & Full-week
calendar and excess-return expert \\
\hline
Meta/hybrid fallback & \code{stage2\_meta\_portfolio\_ensemble.py},
\code{stage2\_hybrid\_gate.py},
\code{stage2\_regularized\_consensus.py},
\code{lstm\_rank\_weight.py} & Fallback chain for non-cached routes \\
\hline
Broad defensive tilt & Internal final-route function & Broad low-risk
route for mild positive flat tape \\
\end{longtable}

The gate uses only as-of observable variables:

\begin{longtable}[]{@{}|
  >{\RaggedRight\arraybackslash}m{(\linewidth - 4\tabcolsep - 3\arrayrulewidth) * \real{0.5000}}|
  >{\RaggedRight\arraybackslash}m{(\linewidth - 4\tabcolsep - 3\arrayrulewidth) * \real{0.5000}}|@{}}
\hline\rowcolor{tableHead}
\begin{minipage}[c]{\linewidth}\centering\bfseries
Gate variable
\end{minipage} & \begin{minipage}[c]{\linewidth}\centering\bfseries
Interpretation
\end{minipage} \\
\hline
\endhead
\hline
\endlastfoot
\code{idx\_ret\_5d} & Recent CSI500 index return \\
\hline
\code{idx\_ret\_20d} & Medium-term index trend \\
\hline
\code{breadth\_ret\_5d\_pos} & Fraction of stocks with positive recent
five-day return \\
\hline
\code{median\_ret\_5d} & Median stock return over the recent five-day
lookback \\
\end{longtable}

The final as-of state was:

\begin{longtable}[]{@{}|>{\RaggedRight\arraybackslash}m{0.58\linewidth}|>{\Centering\arraybackslash}m{0.32\linewidth}|@{}}
\hline\rowcolor{tableHead}
\multicolumn{1}{|c|}{\textbf{Field}} & \multicolumn{1}{c|}{\textbf{Value}} \\
\hline
\endhead
\hline
\endlastfoot
As-of date & 2026-05-08 \\
\hline
Selected route & \code{baseline\_xgb} \\
\hline
Guard reason & \code{baseline\_guard\_overheated\_high\_breadth} \\
\hline
\code{idx\_ret\_5d} & +5.975\% \\
\hline
\code{idx\_ret\_20d} & +14.833\% \\
\hline
\code{breadth\_ret\_5d\_pos} & 70.741\% \\
\hline
Selected stocks & 30 \\
\hline
Max weight & 6.452\% \\
\hline
Effective names & 22.87 \\
\hline
Route validation rank IC & 0.0783 \\
\end{longtable}

This does not mean the final project is simply the unchanged baseline.
It means the full gated model chose the baseline fallback for this as-of
date because historical similar regimes favored a defensive branch over
aggressive weekly consensus.

\subsection{2.8 Portfolio Construction And
Weighting}\label{portfolio-construction-and-weighting}

The portfolio layer was tuned separately from the prediction layer. This
is important because a model can rank stocks reasonably but still lose
excess return if the final portfolio is too diluted, too concentrated,
or too equal-weighted. The final portfolio is not equal-weighted.

\begin{longtable}[]{@{}|>{\RaggedRight\arraybackslash}m{0.58\linewidth}|>{\Centering\arraybackslash}m{0.32\linewidth}|@{}}
\hline\rowcolor{tableHead}
\multicolumn{1}{|c|}{\textbf{Metric}} & \multicolumn{1}{c|}{\textbf{Final value}} \\
\hline
\endhead
\hline
\endlastfoot
Holdings & 30 \\
\hline
Weight sum & 1.0000 \\
\hline
Max weight & 6.452\% \\
\hline
Min weight & 0.215\% \\
\hline
Effective number of names & 22.87 \\
\hline
Top-5 weight & 30.11\% \\
\hline
Top-10 weight & 54.84\% \\
\end{longtable}

The final top holdings are:

\begin{longtable}[]{@{}|>{\Centering\arraybackslash}m{0.10\linewidth}|>{\Centering\arraybackslash}m{0.55\linewidth}|>{\Centering\arraybackslash}m{0.25\linewidth}|@{}}
\hline\rowcolor{tableHead}
\multicolumn{1}{|c|}{\textbf{Rank}} & \multicolumn{1}{c|}{\textbf{Stock\_code}} & \multicolumn{1}{c|}{\textbf{Weight}} \\
\hline
\endhead
\hline
\endlastfoot
1 & 688615 & 6.452\% \\
\hline
2 & 002624 & 6.237\% \\
\hline
3 & 300857 & 6.022\% \\
\hline
4 & 300676 & 5.806\% \\
\hline
5 & 688017 & 5.591\% \\
\hline
6 & 603341 & 5.376\% \\
\hline
7 & 002131 & 5.161\% \\
\hline
8 & 600536 & 4.946\% \\
\hline
9 & 300454 & 4.731\% \\
\hline
10 & 601615 & 4.516\% \\
\end{longtable}

The table below is a \textbf{Portfolio-shape sanity check}, not the main
12-window model comparison. Its purpose was narrower: after the final
2026-05-08 regime gate selected the \code{baseline\_xgb} fallback, we
needed to decide whether that fallback branch should hold the minimum 30
names or dilute into more stocks. To avoid tuning on the final unknown
window, I only used historical windows where the same final gate also
routed into the baseline-style fallback, then rescored the same fallback
idea with \code{top\_k} set to 30, 35, 40, 50, and 60.

The two historical windows used for this check were:

\begin{longtable}[]{@{}|>{\Centering\arraybackslash}m{0.28\linewidth}|>{\Centering\arraybackslash}m{0.34\linewidth}|>{\RaggedRight\arraybackslash}m{0.28\linewidth}|@{}}
\hline\rowcolor{tableHead}
\multicolumn{1}{|c|}{\textbf{As\_of}} & \multicolumn{1}{c|}{\textbf{Evaluation window}} & \multicolumn{1}{c|}{\textbf{Gate reason}} \\
\hline
\endhead
\hline
\endlastfoot
2026-01-23 & 2026-01-26 to 2026-01-30 & Overheated high breadth \\
\hline
2026-03-20 & 2026-03-23 to 2026-03-27 & Severe broad selloff \\
\end{longtable}

Across only those two baseline-routed sanity-check windows, top30 had
the best mean and best minimum:

\begin{longtable}[]{@{}|>{\Centering\arraybackslash}m{0.29\linewidth}|>{\Centering\arraybackslash}m{0.29\linewidth}|>{\Centering\arraybackslash}m{0.29\linewidth}|@{}}
\hline\rowcolor{tableHead}
\multicolumn{1}{|c|}{\textbf{Top\_k}} & \multicolumn{1}{c|}{\textbf{Mean excess}} & \multicolumn{1}{c|}{\textbf{Min excess}} \\
\hline
\endhead
\hline
\endlastfoot
30 & +1.993\% & +1.141\% \\
\hline
35 & +1.831\% & +0.616\% \\
\hline
40 & +1.626\% & +0.153\% \\
\hline
50 & +1.268\% & -0.380\% \\
\hline
60 & +1.072\% & -0.551\% \\
\end{longtable}

The causal use of this table is therefore limited but clear: it did not
decide the overall model, and it was not used to claim that 30 names is
globally optimal for all regimes. It only supported the final
fallback-branch portfolio shape. In the closest available
baseline-routed regimes, larger portfolios diluted the XGBoost rank
signal and lowered the floor. Since the final top30 shape stayed below
the 10\% cap and had 22.87 effective names, I kept top30 for the final
overheated high-breadth fallback.

\section{3. Results}\label{results}

This section reports how the final portfolio generator performed on
held-out complete workweek windows relative to the original XGBoost
baseline. The construction of the train / validation / test split and
the no-leakage checks are documented separately in Section 5, so this
section focuses on the empirical outcome.

\subsection{3.1 Held-Out Performance Relative To
Baseline}\label{held-out-performance-relative-to-baseline}

The final route was compared with the original baseline XGBoost on the
same 12 windows.

\begin{longtable}[]{@{}|
  >{\RaggedRight\arraybackslash}m{(\linewidth - 16\tabcolsep - 9\arrayrulewidth) * \real{0.0968}}|
  >{\Centering\arraybackslash}m{(\linewidth - 16\tabcolsep - 9\arrayrulewidth) * \real{0.1290}}|
  >{\Centering\arraybackslash}m{(\linewidth - 16\tabcolsep - 9\arrayrulewidth) * \real{0.1290}}|
  >{\Centering\arraybackslash}m{(\linewidth - 16\tabcolsep - 9\arrayrulewidth) * \real{0.1290}}|
  >{\Centering\arraybackslash}m{(\linewidth - 16\tabcolsep - 9\arrayrulewidth) * \real{0.1290}}|
  >{\Centering\arraybackslash}m{(\linewidth - 16\tabcolsep - 9\arrayrulewidth) * \real{0.1290}}|
  >{\Centering\arraybackslash}m{(\linewidth - 16\tabcolsep - 9\arrayrulewidth) * \real{0.1290}}|
  >{\Centering\arraybackslash}m{(\linewidth - 16\tabcolsep - 9\arrayrulewidth) * \real{0.1290}}|@{}}
\hline\rowcolor{tableHead}
\begin{minipage}[c]{\linewidth}\centering\bfseries
Model
\end{minipage} & \begin{minipage}[c]{\linewidth}\centering\bfseries
Mean excess
\end{minipage} & \begin{minipage}[c]{\linewidth}\centering\bfseries
Sum excess
\end{minipage} & \begin{minipage}[c]{\linewidth}\centering\bfseries
Median excess
\end{minipage} & \begin{minipage}[c]{\linewidth}\centering\bfseries
Min excess
\end{minipage} & \begin{minipage}[c]{\linewidth}\centering\bfseries
Max excess
\end{minipage} & \begin{minipage}[c]{\linewidth}\centering\bfseries
Count
\end{minipage} & \begin{minipage}[c]{\linewidth}\centering\bfseries
Negative windows
\end{minipage} \\
\hline
\endhead
\hline
\endlastfoot
Final regime-gated ensemble & +4.001\% & +48.006\% & +3.001\% & +0.907\%
& +12.308\% & 12 & 0 \\
\hline
Original XGBoost baseline & +0.600\% & +7.194\% & +0.264\% & -1.221\% &
+2.916\% & 12 & 5 \\
\end{longtable}

The final route improved mean excess by +3.401 percentage points per
held-out window and removed all negative excess windows. It beat the
original baseline in all 12 test windows.

\subsection{3.2 Window-By-Window Held-Out
Results}\label{window-by-window-held-out-results}

\begin{longtable}[]{@{}|>{\Centering\arraybackslash}m{0.25\linewidth}|>{\Centering\arraybackslash}m{0.15\linewidth}|>{\Centering\arraybackslash}m{0.15\linewidth}|>{\Centering\arraybackslash}m{0.15\linewidth}|>{\Centering\arraybackslash}m{0.20\linewidth}|@{}}
\hline\rowcolor{tableHead}
\multicolumn{1}{|c|}{\textbf{Window}} & \multicolumn{1}{c|}{\textbf{Baseline excess}} & \multicolumn{1}{c|}{\textbf{Final excess}} & \multicolumn{1}{c|}{\textbf{Difference}} & \multicolumn{1}{c|}{\textbf{Winner}} \\
\hline
\endhead
\hline
\endlastfoot
2026-01-12 to 2026-01-16 & +2.227\% & +2.948\% & +0.721\% & Final \\
\hline
2026-01-19 to 2026-01-23 & -0.581\% & +6.633\% & +7.214\% & Final \\
\hline
2026-01-26 to 2026-01-30 & -0.380\% & +1.141\% & +1.521\% & Final \\
\hline
2026-02-02 to 2026-02-06 & -0.228\% & +2.741\% & +2.969\% & Final \\
\hline
2026-02-09 to 2026-02-13 & +0.392\% & +6.872\% & +6.480\% & Final \\
\hline
2026-03-02 to 2026-03-06 & -1.221\% & +3.253\% & +4.474\% & Final \\
\hline
2026-03-09 to 2026-03-13 & +0.136\% & +12.308\% & +12.172\% & Final \\
\hline
2026-03-16 to 2026-03-20 & +1.325\% & +2.211\% & +0.886\% & Final \\
\hline
2026-03-23 to 2026-03-27 & +2.916\% & +3.098\% & +0.182\% & Final \\
\hline
2026-03-30 to 2026-04-03 & -0.392\% & +0.907\% & +1.299\% & Final \\
\hline
2026-04-13 to 2026-04-17 & +1.830\% & +2.841\% & +1.011\% & Final \\
\hline
2026-04-20 to 2026-04-24 & +1.170\% & +3.054\% & +1.884\% & Final \\
\end{longtable}

\subsection{3.3 Route Diagnostics For The Reported
Results}\label{route-diagnostics-for-the-reported-results}

The final model is adaptive. Different windows selected different route
types based on as-of regime:

\begin{longtable}[]{@{}|
  >{\Centering\arraybackslash}m{(\linewidth - 14\tabcolsep - 8\arrayrulewidth) * \real{0.1250}}|
  >{\Centering\arraybackslash}m{(\linewidth - 14\tabcolsep - 8\arrayrulewidth) * \real{0.1250}}|
  >{\RaggedRight\arraybackslash}m{(\linewidth - 14\tabcolsep - 8\arrayrulewidth) * \real{0.1250}}|
  >{\RaggedRight\arraybackslash}m{(\linewidth - 14\tabcolsep - 8\arrayrulewidth) * \real{0.1250}}|
  >{\Centering\arraybackslash}m{(\linewidth - 14\tabcolsep - 8\arrayrulewidth) * \real{0.1667}}|
  >{\Centering\arraybackslash}m{(\linewidth - 14\tabcolsep - 8\arrayrulewidth) * \real{0.1667}}|
  >{\Centering\arraybackslash}m{(\linewidth - 14\tabcolsep - 8\arrayrulewidth) * \real{0.1667}}|@{}}
\hline\rowcolor{tableHead}
\begin{minipage}[c]{\linewidth}\centering\bfseries
As\_of
\end{minipage} & \begin{minipage}[c]{\linewidth}\centering\bfseries
Evaluation window
\end{minipage} & \begin{minipage}[c]{\linewidth}\centering\bfseries
Selected route
\end{minipage} & \begin{minipage}[c]{\linewidth}\centering\bfseries
Guard reason
\end{minipage} & \begin{minipage}[c]{\linewidth}\centering\bfseries
\code{idx\_ret\_5d}
\end{minipage} & \begin{minipage}[c]{\linewidth}\centering\bfseries
\code{idx\_ret\_20d}
\end{minipage} & \begin{minipage}[c]{\linewidth}\centering\bfseries
Breadth
\end{minipage} \\
\hline
\endhead
\hline
\endlastfoot
2026-01-09 & 2026-01-12 to 2026-01-16 & Weekly cycle tree & Extreme
broad rebound & +7.92\% & +12.59\% & 88.78\% \\
\hline
2026-01-16 & 2026-01-19 to 2026-01-23 & Weekly consensus & Weekly
consensus allowed & +2.18\% & +15.34\% & 48.10\% \\
\hline
2026-01-23 & 2026-01-26 to 2026-01-30 & Baseline XGBoost & Overheated
high breadth & +4.34\% & +16.84\% & 79.76\% \\
\hline
2026-01-30 & 2026-02-02 to 2026-02-06 & Weekly consensus & Weekly
consensus allowed & -2.56\% & +12.12\% & 27.25\% \\
\hline
2026-02-06 & 2026-02-09 to 2026-02-13 & Weekly consensus & Weekly
consensus allowed & -2.68\% & +1.11\% & 29.86\% \\
\hline
2026-02-27 & 2026-03-02 to 2026-03-06 & Broad defensive tilt & Mild
positive flat tape & +2.79\% & +3.23\% & 64.13\% \\
\hline
2026-03-06 & 2026-03-09 to 2026-03-13 & Weekly consensus & Weekly
consensus allowed & -3.44\% & -1.85\% & 24.65\% \\
\hline
2026-03-13 & 2026-03-16 to 2026-03-20 & Weekly alpha current & Weak flat
positive medium trend & -1.44\% & +1.15\% & 40.88\% \\
\hline
2026-03-20 & 2026-03-23 to 2026-03-27 & Baseline XGBoost & Severe broad
selloff & -5.82\% & -7.88\% & 11.82\% \\
\hline
2026-03-27 & 2026-03-30 to 2026-04-03 & Weekly cycle tree & Long selloff
stabilization & -0.29\% & -10.64\% & 38.68\% \\
\hline
2026-04-10 & 2026-04-13 to 2026-04-17 & Weekly alpha auto & Overheated
negative medium trend & +4.79\% & -4.67\% & 73.15\% \\
\hline
2026-04-17 & 2026-04-20 to 2026-04-24 & Weekly consensus & Weekly
consensus allowed & +3.07\% & +4.28\% & 64.73\% \\
\end{longtable}

This table is important because it shows that the final result is not
produced by selecting the best route after seeing future returns. The
route is chosen from as-of observable market state.

\subsection{3.4 IC As A Secondary Result
Diagnostic}\label{ic-as-a-secondary-result-diagnostic}

In addition to excess return, we used rank-style IC checks as a sanity
signal:

\begin{longtable}[]{@{}|
  >{\RaggedRight\arraybackslash}m{(\linewidth - 6\tabcolsep - 4\arrayrulewidth) * \real{0.3333}}|
  >{\RaggedRight\arraybackslash}m{(\linewidth - 6\tabcolsep - 4\arrayrulewidth) * \real{0.3333}}|
  >{\RaggedRight\arraybackslash}m{(\linewidth - 6\tabcolsep - 4\arrayrulewidth) * \real{0.3333}}|@{}}
\hline\rowcolor{tableHead}
\begin{minipage}[c]{\linewidth}\centering\bfseries
IC type
\end{minipage} & \begin{minipage}[c]{\linewidth}\centering\bfseries
Definition
\end{minipage} & \begin{minipage}[c]{\linewidth}\centering\bfseries
Use
\end{minipage} \\
\hline
\endhead
\hline
\endlastfoot
Validation rank IC & Daily Spearman correlation between validation
target and model prediction & Internal model sanity check during
training \\
\hline
Allocation IC, all universe & Spearman correlation between submitted
weights, including zeros for unheld names, and realized future stock
returns & Checks whether portfolio allocation direction aligns with
realized cross-sectional returns \\
\hline
Allocation IC, selected only & Spearman correlation between
selected-stock weights and realized returns inside the portfolio &
Checks whether heavier selected names outperform lighter selected
names \\
\end{longtable}

IC was not used as the only selection metric. The official objective is
excess return, so IC was used to diagnose ranking quality and avoid
misleading high-return accidents.

\begin{longtable}[]{@{}|
  >{\RaggedRight\arraybackslash}m{(\linewidth - 14\tabcolsep - 8\arrayrulewidth) * \real{0.1111}}|
  >{\Centering\arraybackslash}m{(\linewidth - 14\tabcolsep - 8\arrayrulewidth) * \real{0.1481}}|
  >{\Centering\arraybackslash}m{(\linewidth - 14\tabcolsep - 8\arrayrulewidth) * \real{0.1481}}|
  >{\Centering\arraybackslash}m{(\linewidth - 14\tabcolsep - 8\arrayrulewidth) * \real{0.1481}}|
  >{\Centering\arraybackslash}m{(\linewidth - 14\tabcolsep - 8\arrayrulewidth) * \real{0.1481}}|
  >{\Centering\arraybackslash}m{(\linewidth - 14\tabcolsep - 8\arrayrulewidth) * \real{0.1481}}|
  >{\Centering\arraybackslash}m{(\linewidth - 14\tabcolsep - 8\arrayrulewidth) * \real{0.1481}}|@{}}
\hline\rowcolor{tableHead}
\begin{minipage}[c]{\linewidth}\centering\bfseries
Model
\end{minipage} & \begin{minipage}[c]{\linewidth}\centering\bfseries
Mean excess
\end{minipage} & \begin{minipage}[c]{\linewidth}\centering\bfseries
Min excess
\end{minipage} & \begin{minipage}[c]{\linewidth}\centering\bfseries
Negative windows
\end{minipage} & \begin{minipage}[c]{\linewidth}\centering\bfseries
Mean all-universe allocation IC
\end{minipage} & \begin{minipage}[c]{\linewidth}\centering\bfseries
Median all-universe allocation IC
\end{minipage} & \begin{minipage}[c]{\linewidth}\centering\bfseries
Mean selected-only allocation IC
\end{minipage} \\
\hline
\endhead
\hline
\endlastfoot
Final regime-gated ensemble & +4.001\% & +0.907\% & 0 & 0.0879 & 0.0671
& 0.3055 \\
\hline
Original XGBoost baseline & +0.600\% & -1.221\% & 5 & 0.0025 & -0.0310 &
-0.0010 \\
\end{longtable}

The final route had both stronger excess return and better allocation
IC. This supports the interpretation that the improvement was not only
from market beta or a single lucky window; the final portfolio weights
were more aligned with future cross-sectional winners.

\section{4. Analysis}\label{analysis}

\subsection{4.1 What Worked}\label{what-worked}

The most successful design choices were:

\begin{longtable}[]{@{}|
  >{\RaggedRight\arraybackslash}m{(\linewidth - 4\tabcolsep - 3\arrayrulewidth) * \real{0.5000}}|
  >{\RaggedRight\arraybackslash}m{(\linewidth - 4\tabcolsep - 3\arrayrulewidth) * \real{0.5000}}|@{}}
\hline\rowcolor{tableHead}
\begin{minipage}[c]{\linewidth}\centering\bfseries
Choice
\end{minipage} & \begin{minipage}[c]{\linewidth}\centering\bfseries
Why it worked
\end{minipage} \\
\hline
\endhead
\hline
\endlastfoot
Separating prediction from portfolio allocation & Prevented unstable
synthetic weight labels from dominating training \\
\hline
Compact core fallback & Reduced overfitting and kept a stable defensive
branch \\
\hline
Route-specific richer features & Allowed weekly alpha features without
contaminating all regimes \\
\hline
As-of regime gate & Avoided forcing one model onto incompatible market
states \\
\hline
Robust losses and regularization & Reduced sensitivity to outlier future
returns \\
\hline
Complete full-week validation & Aligned self-test with the Stage2
target \\
\hline
Leakage audits & Protected against accidentally using future information
or cached results \\
\end{longtable}

The final route's biggest improvement came from not trying to make one
model win every regime. The official-style XGBoost baseline is weak on
average, but it is useful in overheated or selloff regimes. Weekly alpha
and weekly cycle routes can be stronger when the market state is more
normal. The gate turns this into a practical mixture-of-experts system.

\subsection{4.2 What Did Not Work}\label{what-did-not-work}

Several ideas looked promising but were not promoted:

\begin{longtable}[]{@{}|
  >{\RaggedRight\arraybackslash}m{(\linewidth - 4\tabcolsep - 3\arrayrulewidth) * \real{0.5000}}|
  >{\RaggedRight\arraybackslash}m{(\linewidth - 4\tabcolsep - 3\arrayrulewidth) * \real{0.5000}}|@{}}
\hline\rowcolor{tableHead}
\begin{minipage}[c]{\linewidth}\centering\bfseries
Failed or limited idea
\end{minipage} & \begin{minipage}[c]{\linewidth}\centering\bfseries
Reason
\end{minipage} \\
\hline
\endhead
\hline
\endlastfoot
Direct rank/weight learning & Labels were synthetic and unstable;
portfolio weights sometimes became too equal-like \\
\hline
Standalone LSTM & Strong local windows, weak horizon transfer and regime
stability \\
\hline
Transformer-style deep models & Data size and validation instability
made them too risky for final \\
\hline
Noisy external data & Strict cleaning did not improve mean and floor
together \\
\hline
Graph/peer factors & Group influence signals were not robust enough in
held-out windows \\
\hline
Blind feature expansion & Improved single windows but hurt minimum
excess \\
\hline
Overly broad top-k portfolios & Diluted rank signal in the final
fallback branch sanity check \\
\hline
Old generic five-day validation & Overstated performance by mixing
weekend/holiday effects \\
\end{longtable}

The Phase 1 result was also a useful warning. The Stage1 model had a
reasonable local backtest but underperformed the class average in the
official 2026-05-06 to 2026-05-08 window. That pushed the Stage2 design
toward simpler, better-regularized models and stricter self-tests.

\subsection{4.3 Why The Final Route Uses Baseline XGBoost On
2026-05-08}\label{why-the-final-route-uses-baseline-xgboost-on-2026-05-08}

This was a point of confusion during development. The final route
selected \code{baseline\_xgb} for 2026-05-08, but the final model is
not merely the original baseline. It is a gated ensemble that sometimes
routes to baseline when market conditions match a defensive regime.

For 2026-05-08:

\begin{longtable}[]{@{}|>{\RaggedRight\arraybackslash}m{0.42\linewidth}|>{\RaggedRight\arraybackslash}m{0.48\linewidth}|@{}}
\hline\rowcolor{tableHead}
\multicolumn{1}{|c|}{\textbf{Signal}} & \multicolumn{1}{c|}{\textbf{Interpretation}} \\
\hline
\endhead
\hline
\endlastfoot
\code{idx\_ret\_5d\ =\ +5.975\%} & Very strong short-term index
rebound \\
\hline
\code{idx\_ret\_20d\ =\ +14.833\%} & Extended medium-term strength \\
\hline
\code{breadth\_ret\_5d\_pos\ =\ 70.741\%} & Broad participation \\
\hline
Guard reason & Overheated high-breadth rebound \\
\end{longtable}

In this state, aggressive weekly consensus can chase already crowded
winners. Historical similar states favored a more conservative XGBoost
fallback with 30 rank-weighted names.

\subsection{4.4 Remaining Limitations}\label{remaining-limitations}

The final route is stronger than the official baseline in our self-test,
but it is not perfect.

\begin{longtable}[]{@{}|
  >{\RaggedRight\arraybackslash}m{(\linewidth - 4\tabcolsep - 3\arrayrulewidth) * \real{0.5000}}|
  >{\RaggedRight\arraybackslash}m{(\linewidth - 4\tabcolsep - 3\arrayrulewidth) * \real{0.5000}}|@{}}
\hline\rowcolor{tableHead}
\begin{minipage}[c]{\linewidth}\centering\bfseries
Limitation
\end{minipage} & \begin{minipage}[c]{\linewidth}\centering\bfseries
Consequence
\end{minipage} \\
\hline
\endhead
\hline
\endlastfoot
Only official OHLCV-style data in final route & No fundamental or
industry structure in final submission \\
\hline
Regime gate is rule-based & May miss nuanced transitions \\
\hline
Self-test windows are limited & 12 full-week windows are useful but
still a small sample \\
\hline
Final as-of selected fallback & Upside may be lower than a more
aggressive route if rebound continues \\
\hline
No true industry-neutral optimizer & Sector concentration risk may
remain \\
\end{longtable}

Future improvements should focus on clean industry/fundamental features,
stronger walk-forward validation, and a more formal meta-learner for
route selection. Any additional data should pass the same strict
lagging, cleaning, correlation, and ablation standards used here.

\section{5. Self-Test}\label{self-test}

This section is written to match the self-test requirement directly. It
explains how the provided data were split into training, validation, and
test sets, why the split is time-safe, how leakage was checked, and what
performance the final model achieved on the held-out test set relative
to the baseline.

\subsection{5.1 Chronological Train / Validation / Test
Split}\label{chronological-train-validation-test-split}

The self-test uses a walk-forward chronological split, not a random
split. This is necessary for financial data because random splitting
would mix future market regimes into training and produce look-ahead
leakage.

The provided data are split separately for each historical as-of date:

\begin{longtable}[]{@{}|
  >{\RaggedRight\arraybackslash}m{(\linewidth - 6\tabcolsep - 4\arrayrulewidth) * \real{0.2200}}|
  >{\RaggedRight\arraybackslash}m{(\linewidth - 6\tabcolsep - 4\arrayrulewidth) * \real{0.5200}}|
  >{\RaggedRight\arraybackslash}m{(\linewidth - 6\tabcolsep - 4\arrayrulewidth) * \real{0.2600}}|@{}}
\hline\rowcolor{tableHead}
\begin{minipage}[c]{\linewidth}\centering\bfseries
Split
\end{minipage} & \begin{minipage}[c]{\linewidth}\centering\bfseries
Definition
\end{minipage} & \begin{minipage}[c]{\linewidth}\centering\bfseries
Purpose
\end{minipage} \\
\hline
\endhead
\hline
\endlastfoot
Training set & Historical supervised rows whose five-day targets are
fully observable before the validation/embargo boundary & Fit model
parameters \\
\hline
Embargo & Five trading days removed between train and validation &
Prevent overlapping five-day target windows \\
\hline
Validation set & The last 10 eligible trading days before the as-of
prediction point, after embargo & Early stopping, validation IC, route
sanity checks \\
\hline
Test set & The future five trading days after as-of, never visible
during feature construction or training & Held-out scoring with
\code{score\_submission.py} \\
\end{longtable}

For every tested as-of date:

\begin{Shaded}
\begin{Highlighting}[]
\NormalTok{1. truncate stock and index data to date \textless{}= as\_of}
\NormalTok{2. build features on the truncated panel}
\NormalTok{3. remove the last 5 trading days from supervised training target availability}
\NormalTok{4. split the remaining supervised data into train / embargo / validation}
\NormalTok{5. train the model route}
\NormalTok{6. generate portfolio at as\_of}
\NormalTok{7. score only on future dates after as\_of}
\end{Highlighting}
\end{Shaded}

The official-style XGBoost branch uses the following concrete split:

\begin{longtable}[]{@{}|>{\RaggedRight\arraybackslash}m{0.42\linewidth}|>{\RaggedRight\arraybackslash}m{0.48\linewidth}|@{}}
\hline\rowcolor{tableHead}
\multicolumn{1}{|c|}{\textbf{Split component}} & \multicolumn{1}{c|}{\textbf{Definition}} \\
\hline
\endhead
\hline
\endlastfoot
Training target cutoff & \code{as\_of\ -\ 5\ trading\ days} \\
\hline
Training set & All eligible rows up to \code{train\_end} \\
\hline
Embargo & 5 trading days discarded between train and validation \\
\hline
Validation set & Last 10 trading days after embargo \\
\hline
Test set & Future five trading days after as-of \\
\end{longtable}

Each portfolio is generated using only data available before its own
test window. The route gate also uses only as-of market state, not
realized future returns.

\subsection{5.2 Validation Design And Window
Correction}\label{validation-design-and-window-correction}

Early in the project, many ``five-day'' tests used the next five trading
days after an arbitrary as-of date. That is a valid generic
holding-period test, but it often crosses weekends or holidays. After
auditing the intended Stage2 target, we changed the main self-test to
complete Monday-Friday workweeks.

\begin{longtable}[]{@{}|
  >{\RaggedRight\arraybackslash}m{(\linewidth - 6\tabcolsep - 4\arrayrulewidth) * \real{0.3000}}|
  >{\Centering\arraybackslash}m{(\linewidth - 6\tabcolsep - 4\arrayrulewidth) * \real{0.4000}}|
  >{\RaggedRight\arraybackslash}m{(\linewidth - 6\tabcolsep - 4\arrayrulewidth) * \real{0.3000}}|@{}}
\hline\rowcolor{tableHead}
\begin{minipage}[c]{\linewidth}\centering\bfseries
Old validation set
\end{minipage} & \begin{minipage}[c]{\linewidth}\centering\bfseries
Complete Mon-Fri windows
\end{minipage} & \begin{minipage}[c]{\linewidth}\centering\bfseries
Issue
\end{minipage} \\
\hline
\endhead
\hline
\endlastfoot
Old 12-window setup & 2 / 12 & Most windows crossed weekend or holiday
gaps \\
\hline
Old 9-window setup & 0 / 9 & None were strict workweeks \\
\end{longtable}

This correction invalidated an earlier optimistic high-water mark around
+8.7\% mean excess. It was useful research evidence but not safe as the
final Stage2 standard. The final model-selection criterion therefore
prioritized the complete-workweek test set, while also checking that the
route did not rely on one lucky window.

\subsection{5.3 Held-Out Test Set And
Performance}\label{held-out-test-set-and-performance}

The final 12 held-out complete workweek windows were:

\begin{longtable}[]{@{}|>{\Centering\arraybackslash}m{0.42\linewidth}|>{\Centering\arraybackslash}m{0.48\linewidth}|@{}}
\hline\rowcolor{tableHead}
\multicolumn{1}{|c|}{\textbf{As\_of}} & \multicolumn{1}{c|}{\textbf{Evaluation window}} \\
\hline
\endhead
\hline
\endlastfoot
2026-01-09 & 2026-01-12 to 2026-01-16 \\
\hline
2026-01-16 & 2026-01-19 to 2026-01-23 \\
\hline
2026-01-23 & 2026-01-26 to 2026-01-30 \\
\hline
2026-01-30 & 2026-02-02 to 2026-02-06 \\
\hline
2026-02-06 & 2026-02-09 to 2026-02-13 \\
\hline
2026-02-27 & 2026-03-02 to 2026-03-06 \\
\hline
2026-03-06 & 2026-03-09 to 2026-03-13 \\
\hline
2026-03-13 & 2026-03-16 to 2026-03-20 \\
\hline
2026-03-20 & 2026-03-23 to 2026-03-27 \\
\hline
2026-03-27 & 2026-03-30 to 2026-04-03 \\
\hline
2026-04-10 & 2026-04-13 to 2026-04-17 \\
\hline
2026-04-17 & 2026-04-20 to 2026-04-24 \\
\end{longtable}

The final model's performance on this held-out test set was:

\begin{longtable}[]{@{}|
  >{\RaggedRight\arraybackslash}m{(\linewidth - 10\tabcolsep - 6\arrayrulewidth) * \real{0.1900}}|
  >{\Centering\arraybackslash}m{(\linewidth - 10\tabcolsep - 6\arrayrulewidth) * \real{0.1620}}|
  >{\Centering\arraybackslash}m{(\linewidth - 10\tabcolsep - 6\arrayrulewidth) * \real{0.1620}}|
  >{\Centering\arraybackslash}m{(\linewidth - 10\tabcolsep - 6\arrayrulewidth) * \real{0.1620}}|
  >{\Centering\arraybackslash}m{(\linewidth - 10\tabcolsep - 6\arrayrulewidth) * \real{0.1620}}|
  >{\Centering\arraybackslash}m{(\linewidth - 10\tabcolsep - 6\arrayrulewidth) * \real{0.1620}}|@{}}
\hline\rowcolor{tableHead}
\begin{minipage}[c]{\linewidth}\centering\bfseries
Model
\end{minipage} & \begin{minipage}[c]{\linewidth}\centering\bfseries
Mean test excess
\end{minipage} & \begin{minipage}[c]{\linewidth}\centering\bfseries
Median test excess
\end{minipage} & \begin{minipage}[c]{\linewidth}\centering\bfseries
Min test excess
\end{minipage} & \begin{minipage}[c]{\linewidth}\centering\bfseries
Negative test windows
\end{minipage} & \begin{minipage}[c]{\linewidth}\centering\bfseries
Wins versus baseline
\end{minipage} \\
\hline
\endhead
\hline
\endlastfoot
Final regime-gated ensemble & +4.001\% & +3.001\% & +0.907\% & 0 / 12
& 12 / 12 \\
\hline
Original XGBoost baseline & +0.600\% & +0.264\% & -1.221\% & 5 / 12
& 0 / 12 \\
\end{longtable}

This satisfies the self-test requirement that the reported test
performance exceed the provided baseline. The final model improved mean
test excess by +3.401 percentage points and had no negative excess test
windows, while the baseline had five negative windows. The detailed
window-by-window returns are shown in Section 3.2.

\subsection{5.4 No-Leakage And Cache
Controls}\label{no-leakage-and-cache-controls}

We performed both dynamic and static checks.

\begin{longtable}[]{@{}|
  >{\RaggedRight\arraybackslash}m{(\linewidth - 6\tabcolsep - 4\arrayrulewidth) * \real{0.3333}}|
  >{\RaggedRight\arraybackslash}m{(\linewidth - 6\tabcolsep - 4\arrayrulewidth) * \real{0.3333}}|
  >{\RaggedRight\arraybackslash}m{(\linewidth - 6\tabcolsep - 4\arrayrulewidth) * \real{0.3333}}|@{}}
\hline\rowcolor{tableHead}
\begin{minipage}[c]{\linewidth}\centering\bfseries
Check
\end{minipage} & \begin{minipage}[c]{\linewidth}\centering\bfseries
Purpose
\end{minipage} & \begin{minipage}[c]{\linewidth}\centering\bfseries
Result
\end{minipage} \\
\hline
\endhead
\hline
\endlastfoot
Format validation & Verify stock count, positive weights, sum to one,
max weight cap & Passed \\
\hline
Dynamic truncation audit & Compare full-data generation with physically
truncated \code{date\ \textless{}=\ as\_of} data & Passed \\
\hline
Static code scan & Search active route files for
\code{score\_submission}, report-cache, archive, and leakage-prone
references & Reviewed \\
\hline
Regeneration check & Regenerate final
\code{stage2\_final\_portfolio.csv} from script and compare & Exact
match \\
\hline
py\_compile & Ensure active scripts import cleanly & Passed \\
\end{longtable}

The final dynamic leakage audit was rerun after the last folder cleanup.
It passed on six representative as-of dates chosen to cover the major
gate branches: weekly cycle, broad defensive tilt, weekly alpha, and the
final baseline fallback.

\begin{longtable}[]{@{}|
  >{\Centering\arraybackslash}m{(\linewidth - 12\tabcolsep - 7\arrayrulewidth) * \real{0.1429}}|
  >{\RaggedRight\arraybackslash}m{(\linewidth - 12\tabcolsep - 7\arrayrulewidth) * \real{0.1429}}|
  >{\Centering\arraybackslash}m{(\linewidth - 12\tabcolsep - 7\arrayrulewidth) * \real{0.1905}}|
  >{\Centering\arraybackslash}m{(\linewidth - 12\tabcolsep - 7\arrayrulewidth) * \real{0.1905}}|
  >{\Centering\arraybackslash}m{(\linewidth - 12\tabcolsep - 7\arrayrulewidth) * \real{0.1905}}|
  >{\Centering\arraybackslash}m{(\linewidth - 12\tabcolsep - 7\arrayrulewidth) * \real{0.1429}}|@{}}
\hline\rowcolor{tableHead}
\begin{minipage}[c]{\linewidth}\centering\bfseries
As\_of
\end{minipage} & \begin{minipage}[c]{\linewidth}\centering\bfseries
Model
\end{minipage} & \begin{minipage}[c]{\linewidth}\centering\bfseries
Max absolute weight difference
\end{minipage} & \begin{minipage}[c]{\linewidth}\centering\bfseries
L1 difference
\end{minipage} & \begin{minipage}[c]{\linewidth}\centering\bfseries
Changed names
\end{minipage} & \begin{minipage}[c]{\linewidth}\centering\bfseries
Pass
\end{minipage} \\
\hline
\endhead
\hline
\endlastfoot
2026-01-09 & Final route & 0.0 & 0.0 & 0 & True \\
\hline
2026-02-27 & Final route & 0.0 & 0.0 & 0 & True \\
\hline
2026-03-13 & Final route & 0.0 & 0.0 & 0 & True \\
\hline
2026-03-27 & Final route & 0.0 & 0.0 & 0 & True \\
\hline
2026-04-10 & Final route & 0.0 & 0.0 & 0 & True \\
\hline
2026-05-08 & Final route & 0.0 & 0.0 & 0 & True \\
\end{longtable}

The active route trims stock and index data to
\code{date\ \textless{}=\ as\_of} before feature construction. Target
columns use forward shifts only for supervised training rows, and
prediction rows do not use future targets. Optional cache arguments
exist for speed during research, but the final reproduction path and
leakage audit run with no cache directory; the final route regenerates
live from the official data cache.

\section{6. Reproducibility}\label{reproducibility}

\subsection{6.1 Environment}\label{environment}

\begin{Shaded}
\begin{Highlighting}[]
\ExtensionTok{conda}\NormalTok{ activate mlcomp{-}sp26}
\ExtensionTok{pip}\NormalTok{ install }\AttributeTok{{-}r}\NormalTok{ requirements.txt}
\end{Highlighting}
\end{Shaded}

Main package requirements:

\begin{longtable}[]{@{}|>{\RaggedRight\arraybackslash}m{0.42\linewidth}|>{\Centering\arraybackslash}m{0.48\linewidth}|@{}}
\hline\rowcolor{tableHead}
\multicolumn{1}{|c|}{\textbf{Package}} & \multicolumn{1}{c|}{\textbf{Requirement}} \\
\hline
\endhead
\hline
\endlastfoot
pandas & \textgreater= 2.0 \\
\hline
numpy & \textgreater= 1.24, \textless{} 2.0 \\
\hline
xgboost & \textgreater= 2.0 \\
\hline
lightgbm & \textgreater= 4.6 \\
\hline
torch & \textgreater= 2.0 \\
\hline
scikit-learn & \textgreater= 1.3 \\
\hline
scipy & \textgreater= 1.11 \\
\hline
pyarrow & \textgreater= 14.0 \\
\hline
catboost & \textgreater= 1.2 \\
\end{longtable}

\subsection{6.2 Reproduce Final
Portfolio}\label{reproduce-final-portfolio}

\begin{Shaded}
\begin{Highlighting}[]
\ExtensionTok{python}\NormalTok{ stage2\_baseline\_guard\_ensemble.py }\DataTypeTok{\textbackslash{}}
  \AttributeTok{{-}{-}as{-}of}\NormalTok{ 20260508 }\DataTypeTok{\textbackslash{}}
  \AttributeTok{{-}{-}baseline{-}top{-}k}\NormalTok{ 0 }\DataTypeTok{\textbackslash{}}
  \AttributeTok{{-}{-}out}\NormalTok{ stage2\_final\_portfolio.csv }\DataTypeTok{\textbackslash{}}
  \AttributeTok{{-}{-}meta{-}out}\NormalTok{ stage2\_report/final\_report\_materials/01\_final\_portfolio\_metadata.csv}

\ExtensionTok{python}\NormalTok{ validate\_submission.py stage2\_final\_portfolio.csv}
\end{Highlighting}
\end{Shaded}

\subsection{6.3 Reproduce Full-Week
Self-Test}\label{reproduce-full-week-self-test}

\begin{Shaded}
\begin{Highlighting}[]
\ExtensionTok{python}\NormalTok{ tools/stage2\_validation/stage2\_backtest\_5day.py }\DataTypeTok{\textbackslash{}}
  \AttributeTok{{-}{-}models}\NormalTok{ baseline\_xgb baseline\_guard\_adaptive }\DataTypeTok{\textbackslash{}}
  \AttributeTok{{-}{-}full{-}week{-}only} \DataTypeTok{\textbackslash{}}
  \AttributeTok{{-}{-}windows}\NormalTok{ 12 }\DataTypeTok{\textbackslash{}}
  \AttributeTok{{-}{-}jobs}\NormalTok{ 4 }\DataTypeTok{\textbackslash{}}
  \AttributeTok{{-}{-}out{-}dir}\NormalTok{ stage2\_report/backtests/guard\_route\_v2\_nocache\_fullweek12\_20260510 }\DataTypeTok{\textbackslash{}}
  \AttributeTok{{-}{-}summary{-}out}\NormalTok{ stage2\_report/final\_report\_materials/02\_full\_week\_12\_window\_performance\_summary.csv }\DataTypeTok{\textbackslash{}}
  \AttributeTok{{-}{-}detail{-}out}\NormalTok{ stage2\_report/final\_report\_materials/03\_full\_week\_12\_window\_performance\_detail.csv}
\end{Highlighting}
\end{Shaded}

The \texttt{stage2\_report/backtests/} folder stores the per-window
portfolio CSVs generated by the final 12-window self-test, plus companion
metadata CSVs for the adaptive route. The command above refreshes those
portfolio records together with the retained summary/detail evidence files.

\subsection{6.4 Reproduce Leakage Audit}\label{reproduce-leakage-audit}

\begin{Shaded}
\begin{Highlighting}[]
\ExtensionTok{python}\NormalTok{ tools/stage2\_validation/stage2\_leakage\_audit.py }\DataTypeTok{\textbackslash{}}
  \AttributeTok{{-}{-}as{-}of}\NormalTok{ 20260109 20260227 20260313 20260327 20260410 20260508 }\DataTypeTok{\textbackslash{}}
  \AttributeTok{{-}{-}models}\NormalTok{ baseline\_guard\_adaptive }\DataTypeTok{\textbackslash{}}
  \AttributeTok{{-}{-}out}\NormalTok{ stage2\_report/final\_report\_materials/05\_final\_leakage\_audit\_dynamic.csv }\DataTypeTok{\textbackslash{}}
  \AttributeTok{{-}{-}static{-}out}\NormalTok{ stage2\_report/final\_report\_materials/06\_final\_leakage\_audit\_static\_scan.csv}
\end{Highlighting}
\end{Shaded}

\subsection{6.5 Final Consistency
Checks}\label{final-consistency-checks}

After removing the old draft report and consolidating the project
layout, the final consistency check confirmed:

\begin{longtable}[]{@{}|
  >{\RaggedRight\arraybackslash}m{(\linewidth - 4\tabcolsep - 3\arrayrulewidth) * \real{0.5000}}|
  >{\Centering\arraybackslash}m{(\linewidth - 4\tabcolsep - 3\arrayrulewidth) * \real{0.5000}}|@{}}
\hline\rowcolor{tableHead}
\begin{minipage}[c]{\linewidth}\centering\bfseries
Check
\end{minipage} & \begin{minipage}[c]{\linewidth}\centering\bfseries
Result
\end{minipage} \\
\hline
\endhead
\hline
\endlastfoot
\code{stage2\_final\_portfolio.csv} format validation & Passed \\
\hline
Data cache max date & 2026-05-08 \\
\hline
Final portfolio row count & 30 \\
\hline
Final portfolio weight sum & 1.000000000000 \\
\hline
Final portfolio max weight & 6.4516\% \\
\hline
Final effective number of names & 22.8689 \\
\hline
Regenerated final portfolio versus root CSV & Exact match, max
difference 0 \\
\hline
Performance summary versus detail CSV & Matched \\
\hline
IC summary/detail versus score detail & Matched \\
\hline
Gate route table versus self-test windows & Matched \\
\hline
Dynamic leakage audit & 6 / 6 rows passed \\
\end{longtable}

\subsection{6.6 Report Evidence Files}\label{report-evidence-files}

\begin{longtable}[]{@{}|
  >{\RaggedRight\arraybackslash}m{(\linewidth - 4\tabcolsep - 3\arrayrulewidth) * \real{0.5000}}|
  >{\RaggedRight\arraybackslash}m{(\linewidth - 4\tabcolsep - 3\arrayrulewidth) * \real{0.5000}}|@{}}
\hline\rowcolor{tableHead}
\begin{minipage}[c]{\linewidth}\centering\bfseries
File
\end{minipage} & \begin{minipage}[c]{\linewidth}\centering\bfseries
Purpose
\end{minipage} \\
\hline
\endhead
\hline
\endlastfoot
\code{stage2\_report/final\_report\_materials/01\_final\_portfolio\_metadata.csv}
& Final as-of route metadata and validation IC \\
\hline
\code{stage2\_report/final\_report\_materials/02\_full\_week\_12\_window\_performance\_summary.csv}
& Aggregate self-test performance \\
\hline
\code{stage2\_report/final\_report\_materials/03\_full\_week\_12\_window\_performance\_detail.csv}
& Per-window returns \\
\hline
\code{stage2\_report/final\_report\_materials/04\_final\_topk\_and\_weighting\_check.md}
& Top-k and weight-shape sanity check \\
\hline
\code{stage2\_report/final\_report\_materials/05\_final\_leakage\_audit\_dynamic.csv}
& Dynamic no-leak audit \\
\hline
\code{stage2\_report/final\_report\_materials/06\_final\_leakage\_audit\_static\_scan.csv}
& Static leakage scan \\
\hline
\code{stage2\_report/final\_report\_materials/07\_stage2\_experiment\_history.md}
& Chronological experiment log \\
\hline
\code{stage2\_report/final\_report\_materials/08\_full\_week\_12\_window\_ic\_detail.csv}
& Per-window allocation IC \\
\hline
\code{stage2\_report/final\_report\_materials/09\_full\_week\_12\_window\_ic\_summary.csv}
& Aggregate allocation IC \\
\hline
\code{stage2\_report/final\_report\_materials/10\_full\_week\_gate\_route\_detail.csv}
& Route selected by the regime gate \\
\hline
\code{stage2\_report/backtests/}
& Saved per-window portfolio CSVs and adaptive-route metadata from the final
12-window self-test \\
\hline
\code{stage2\_report/backtests/README.md}
& Explains how to regenerate the per-window portfolio records \\
\end{longtable}

\section{7. Conclusion}\label{conclusion}

The final Stage2 system is a leakage-safe, walk-forward-tested,
regime-gated stock-selection ensemble. The strongest lesson from the
project was that robust portfolio performance came less from adding a
more complex model and more from matching the validation window,
controlling overfitting, separating score prediction from weight
allocation, and routing among experts based only on as-of market state.

The final adaptive route outperformed the original XGBoost baseline on
every one of the 12 strict complete workweek self-test windows, with
+4.001\% mean excess versus +0.600\% for the baseline, and with zero
negative excess windows. The final portfolio for 2026-05-08 is therefore
a defensive branch selected by a broader ensemble, not a simple
unmodified baseline submission.

\end{document}
